% Pursuing Stacks (cotes 134-2 to 134-8): archive cote 134-3, site page 82 (PDF page 83), typescript page 82.
% typescript. Draft compared: ps2.tex:4260-4339.
% AI-transcribed by Claude Opus 5.5 (claude-opus-5-5, adaptive thinking, effort high), Message Batches API, batch ps1 (msgbatch_01Rv5hCaNTtUt4JHx1yqoHPL), 2026-11-12; image review recorded in records/CHECK_DECISIONS_134_3_4.tsv.
\scanpage{134-3/82}
\mspage{82}
% NOTE: faint marks "(82)" right of the page number and "(117)" (unsure) at the top right; they look like show-through or a stamp
(i)  All induced categories $A_{/a}$  (with a in A) are weak test categori\unsure{es}
(ii)  For any aspheric F in $A^{\hand{\wedge}}$, $A_{/F}$ is a weak test category.
(iii)  T\over[typed]{H}{h}ere exists a  ``homotopy interval'' in $A^{\hand{\wedge}}$,  namely an element
\xout{in} I in $A^{\hand{\wedge}}$, \uline{aspheric over the final element} $e_{A^{\hand{\wedge}}}$ (i.e. such that
all products Ixa (for a in A) are  aspheric),  endowed with
two sections $\hand{\delta}_o$ \xout{\illegible{3}} , $\hand{\delta}_1$ over \over{e}{e} = $e_{A^{\hand{\wedge}}}$ i.e. with two subobjects
$e_o,\over[typed]{e}{e}_{\over[typed]{1}{1}}$  \xout{equivalent} isomorphic to e,  such that  Ker$(\hand{\delta}_o,\hand{\delta}_1) = \hand{\emptyset}_{A^{\hand{\wedge}}}$
i.e. $e_o\hand{\cap}e_1 = \hand{\emptyset}_{A^{\hand{\wedge}}}$.
(iv)  The Lawvere element $L_A$ or $L_{A^{\hand{\wedge}}}$ in $A^{\hand{\wedge}}$ (i.e. the \xout{functor} presheaf
a $\hand{\mapsto}$ subobjects of a in $A^{\hand{\wedge}} \hand{\approx}$ cribles of $A_{/a}$ )  is aspheric over
the final element e  (\over[typed]{o}{i}\over[typed]{n}{n}\hand{\unsure{,}} other words, it is a \hand{``}homotopy inter-%
val\hand{''},  when endowed with the two sections $\hand{\delta}_o,\hand{\delta}_1$  corresponding
to the \xout{empty} full and the empty crible in A).
(v)  For any \xout{category} C in (Cat),  the canonical \xout{functor} \inserted[typed]{map in} (Cat)
\xout{where the}
(7) \[ \over{i}{i}_A j_A(C) = A_{/j_A(C)} \hand{\longrightarrow} AxC \]
(with second component (3))  is (not only a weak equivalence,
which definitely isn\hand{'}t enough,  but \over{\unsure{even}}{even}\hand{)} \uline{aspheric} .

By the usual criterion of asphericity for a map in (Cat),  condition
(v)  is equivalent \deleted{with} \inserted{to} the condition  (v') :  For any a in A and any C
in (Cat) with final element,  $A_{/axj_A(C)}$  is aspheric i.e.  $axj_A(C)$  is an
aspheric element in $A^{\hand{\wedge}}$ ;  \deleted{which} \inserted{now this} is clearly equivalent to (i)  (by the
checking-criterion (v) above for weak test categories,  applied to the
categories $A_{/a}$ ). Thus we get the purely formal implications
% Hand-drawn arrows are marked blue; typed objects remain plain.
\diagram{\begin{tikzcd}
(ii) \arrow[r, Rightarrow, "\hand{}",color=blue] & (i) \arrow[r, Rightarrow,color=blue] \arrow[d, Leftrightarrow,color=blue] & (iv) \arrow[r, Rightarrow,color=blue] & (iii) \\
 & (v) & &
\end{tikzcd}}
\deleted{\unsure{,}} ,
% NOTE: the arrows of the diagram are drawn by hand

where (i)\hand{$\Rightarrow$}(iv) is obtained by applying the criterion (v') just recalled
to the case of  C = $\hand{\Delta}_1$ .  As the condition (iii) is clearly stable by
localizing to a category $A_{/F}$  (using  the equivalence  $(A_{/F})^{\hand{\wedge}} \hand{\simeq} A^{\hand{\wedge}}{}_{/F}$ ), to
see that (iii) implies (ii),  we are reduced \inserted{(replacing A by A/F)} to the case when F is the
final element in $A^{\hand{\wedge}}$,  and thus to proving that if A  has a final element,
then (iii) implies that A is a weak test category,  i.e. that the catego-%
ries  $A_{/j_A(C)}$ \hand{,} with C \deleted{in (Cat)} having a final element\hand{,}  are aspheric.
This was done by a simple  homotopy argument  in \xout{the "key result" of pag}
\over{\unsure{o}}{t}\over[typed]{w}{w}o steps.  O\over[typed]{N}{n}e step  (\hand{``}the comparison lemma for homotopy intervals''
on \hand{\S.37})  shows that (iii)\hand{$\Rightarrow$}(iv)  i.e. existence of a homotopy interval
implies that the Lawvere interval is a homotopy interval ,  the other
step  (presented in a more general set-up in the ``key result'' \deleted{on page 61} \inserted{in same \S.31})
proving that (iv)  implies asphericity of the elements $j_A($ C)  (C with
